% =====================================================================
\chapter{Datos de preentrenamiento}
\label{tema:1}
% =====================================================================

El preentrenamiento minimiza la entropía cruzada media sobre un corpus (módulo~1, tema~5), de modo que el modelo aprende la distribución del texto que se le muestra. La composición de ese corpus es, por tanto, una decisión de diseño tan importante como la arquitectura. Este tema describe cómo se obtiene un corpus a partir de la web y formaliza tres de sus etapas: los filtros de calidad, el filtro por perplejidad y la mezcla de fuentes. Se siguen las secciones~2 y~4 de la práctica de datos de CS336~\cite{cs336a4}, la clase sobre datos de CS324~\cite{cs324} y los informes de CCNet~\cite{wenzek}, Gopher~\cite[§3.4 y apéndice~A]{rae}, RefinedWeb~\cite[§3]{refinedweb}, FineWeb~\cite{fineweb} y LLaMA~\cite[§2.1]{llama}. La deduplicación, que también forma parte del proceso, se estudia por separado en el \cref{tema:2}.

\section{Del rastreo web al corpus}

\begin{definicion}[Corpus de preentrenamiento]
Un \emph{corpus de preentrenamiento} es una colección finita de documentos $\mathcal{C}=\{x^{(1)},\dots,x^{(M)}\}$, cada uno una secuencia de tokens. Su \emph{tamaño} es el número total de tokens, $\sum_{m}\abs{x^{(m)}}$. Un \emph{filtro} es un predicado $f$ sobre documentos; aplicarlo a $\mathcal{C}$ produce $\{x\in\mathcal{C}:f(x)=1\}$. Un proceso de preparación es una composición de transformaciones (extracción, normalización) y filtros.
\end{definicion}

La fuente principal es la web. Common Crawl publica volcados periódicos de miles de millones de páginas; FineWeb, por ejemplo, se construyó a partir de 96 volcados y contiene unos $15$ billones de tokens~\cite{fineweb}. La \cref{fig:proceso} resume las etapas habituales. Casi todo el material en bruto se descarta: en RefinedWeb, la identificación del idioma elimina cerca de la mitad de los documentos, y el conjunto del proceso, casi el $90\,\%$~\cite[§3 y figura~2]{refinedweb}.

\begin{figure}[h]
\centering
\begin{tikzpicture}[font=\small,
  caja/.style={draw,rounded corners=2pt,align=center,minimum height=9mm,text width=2.25cm,inner sep=2pt},
  flecha/.style={-{Stealth[length=2mm]},thick}]
  \node[caja] (a) {Rastreo\\(HTML)};
  \node[caja,right=0.45cm of a] (b) {Extracción\\del texto};
  \node[caja,right=0.45cm of b] (c) {Identificación\\del idioma};
  \node[caja,right=0.45cm of c] (d) {Filtros de\\calidad};
  \node[caja,below=0.9cm of d] (e) {Deduplicación\\(\cref{tema:2})};
  \node[caja,left=0.45cm of e] (f) {Datos personales\\y toxicidad};
  \node[caja,left=0.45cm of f] (g) {Mezcla de\\fuentes};
  \node[caja,left=0.45cm of g] (h) {Tokenización\\(módulo~12)};
  \draw[flecha] (a)--(b); \draw[flecha] (b)--(c); \draw[flecha] (c)--(d);
  \draw[flecha] (d)--(e); \draw[flecha] (e)--(f); \draw[flecha] (f)--(g); \draw[flecha] (g)--(h);
\end{tikzpicture}
\caption{Etapas habituales de la preparación de un corpus a partir de la web, según~\cite{cs336a4,refinedweb,fineweb}. El orden concreto varía de un informe a otro.}
\label{fig:proceso}
\end{figure}

\section{Filtros heurísticos de calidad}

Los filtros más sencillos son reglas sobre estadísticas superficiales del documento. Son baratos, deterministas y fáciles de auditar. El informe de Gopher~\cite[apéndice~A]{rae} define un conjunto de ellas; la práctica de CS336~\cite[§2.6]{cs336a4} usa el siguiente subconjunto.

\begin{definicion}[Filtros de Gopher, subconjunto de CS336]\label{def:gopher}
Un documento con $W$ palabras, longitud media de palabra $\bar{\ell}$ (en caracteres) y $\Lambda$ líneas \emph{pasa} el filtro si se cumplen las cuatro condiciones siguientes:
\begin{enumerate}
  \item $50\le W\le 100\,000$;
  \item $3\le\bar{\ell}\le 10$;
  \item como mucho el $30\,\%$ de las $\Lambda$ líneas termina en puntos suspensivos;
  \item al menos el $80\,\%$ de las palabras contiene un carácter alfabético.
\end{enumerate}
\end{definicion}

\begin{ejemplo}\label{ej:gopher}
Un documento de $120$ palabras y $480$ caracteres en palabras, con $10$ líneas, $4$ de las cuales terminan en «...», tiene $\bar{\ell}=480/120=4$, pero una fracción de líneas truncadas de $4/10=0{,}40>0{,}30$. No pasa el filtro: es típico de las páginas de listados con fragmentos de artículos.
\end{ejemplo}

\section{Filtro por perplejidad}

CCNet~\cite{wenzek} puntúa cada página con un modelo de $n$-gramas entrenado sobre Wikipedia y conserva las de perplejidad baja; LLaMA usó este proceso para su parte de Common Crawl~\cite[§2.1]{llama}. La idea es que el texto parecido al de referencia recibe probabilidad alta. El caso de un modelo de unigramas permite analizar exactamente qué mide el filtro.

\begin{definicion}[Perplejidad de un documento]
Sea $q$ una distribución sobre $\V$ con $q(v)>0$ para todo $v$ (un modelo de unigramas de referencia) y sea $x=(x_1,\dots,x_T)$ un documento. La perplejidad de $x$ según $q$ es
\[
\PPL_q(x)=\exp\Bigl(-\frac1T\sum_{t=1}^{T}\ln q(x_t)\Bigr).
\]
La \emph{distribución empírica} de $x$ es $\hat p_x(v)=\abs{\{t:x_t=v\}}/T$.
\end{definicion}

\begin{proposicion}\label{prop:ppl-kl}
Para todo documento $x$,
\[
\ln\PPL_q(x)=H(\hat p_x,q)=H(\hat p_x)+D(\hat p_x\,\|\,q).
\]
\end{proposicion}

\begin{proof}
Agrupando las posiciones por token, $\sum_{t}\ln q(x_t)=\sum_{v}T\hat p_x(v)\ln q(v)$. Dividiendo por $-T$ se obtiene $-\sum_v\hat p_x(v)\ln q(v)=H(\hat p_x,q)$, la entropía cruzada (módulo~1, tema~5). La segunda igualdad es la descomposición $H(p,q)=H(p)+D(p\,\|\,q)$ del módulo~6, que se comprueba sumando y restando $\sum_v\hat p_x(v)\ln\hat p_x(v)$.
\end{proof}

La \cref{prop:ppl-kl} muestra que el filtro premia dos cosas distintas: que el documento se parezca a la referencia (divergencia pequeña) y que el documento sea poco variado (entropía empírica pequeña). La segunda es indeseable.

\begin{corolario}[El filtro premia la repetición]\label{cor:repeticion}
Sea $q_{\max}=\max_v q(v)$. Para todo documento $x$, $\PPL_q(x)\ge 1/q_{\max}$, con igualdad si y solo si todos los tokens de $x$ maximizan $q$.
\end{corolario}

\begin{proof}
Cada factor cumple $-\ln q(x_t)\ge-\ln q_{\max}$, luego la media también, y $\PPL_q(x)\ge e^{-\ln q_{\max}}=1/q_{\max}$. La igualdad en la media exige la igualdad en cada término, es decir, $q(x_t)=q_{\max}$ para todo $t$.
\end{proof}

\begin{ejemplo}\label{ej:ppl-filtro}
Sea $q(\texttt{a})=0{,}5$, $q(\texttt{b})=0{,}3$, $q(\texttt{c})=0{,}2$. El documento «a b a c» tiene $\PPL_q=(0{,}5\cdot 0{,}3\cdot 0{,}5\cdot 0{,}2)^{-1/4}=0{,}015^{-1/4}=2{,}857$. Su distribución empírica es $(1/2,1/4,1/4)$, con $H(\hat p)=1{,}0397$ nats y $D(\hat p\,\|\,q)=0{,}0102$ nats; la suma, $1{,}0499$, es $\ln 2{,}857$, como predice la \cref{prop:ppl-kl}. El documento degenerado «a a a a» obtiene la perplejidad mínima, $1/0{,}5=2$. Por eso los procesos reales combinan el filtro por perplejidad con reglas contra la repetición~\cite[apéndice~A]{rae}.
\end{ejemplo}

\section{Clasificadores de calidad}

Una alternativa es aprender qué es «calidad» a partir de ejemplos. Se toman como positivos documentos de una fuente de referencia (por ejemplo, páginas citadas desde Wikipedia) y como negativos páginas aleatorias de Common Crawl, y se entrena un clasificador lineal~\cite[§2.7]{cs336a4}.

\begin{definicion}[Filtro por clasificador]
Sea $g\colon\mathcal{X}\to[0,1]$ un clasificador que estima la probabilidad de que un documento proceda de la fuente de referencia. El filtro con umbral $\tau$ conserva los documentos con $g(x)\ge\tau$. Subir $\tau$ aumenta la precisión y reduce la cobertura.
\end{definicion}

LLaMA entrenó un clasificador lineal que distingue las páginas usadas como referencia en Wikipedia de páginas aleatorias, y descartó las no clasificadas como referencia~\cite[§2.1]{llama}. FineWeb-Edu sustituye la fuente de referencia por anotaciones de un modelo grande: Llama-3-70B-Instruct puntuó el «valor educativo» de una muestra de páginas de $0$ a $5$, se ajustó una regresión sobre embeddings y se conservaron las páginas con puntuación al menos~$3$, unos $1{,}3$ billones de tokens~\cite[§4]{fineweb}.

\begin{observacion}
«Calidad» es una elección, no una propiedad objetiva del texto. Un clasificador entrenado con Wikipedia como modelo de lo deseable tiende a descartar registros informales, variedades dialectales e idiomas con pocos recursos, y el modelo resultante hereda ese sesgo. Las reglas heurísticas tienen el mismo problema en menor grado: un umbral de longitud media de palabra pensado para el inglés no es neutral para otros idiomas.
\end{observacion}

\section{Mezcla de fuentes y épocas}

El corpus final combina varias fuentes (web, código, libros, artículos científicos, Wikipedia) con pesos fijados por experimentos.

\begin{definicion}[Mezcla y épocas]\label{def:mezcla}
Sean $m$ fuentes con $U_i$ tokens únicos cada una y sean $w_i\ge 0$, $\sum_i w_i=1$, los \emph{pesos de muestreo}. Un entrenamiento de $D$ tokens extrae $w_iD$ tokens de la fuente $i$; su número de \emph{épocas} es
\[
e_i=\frac{w_iD}{U_i}.
\]
Si $e_i>1$, la fuente se \emph{sobremuestrea}: parte de su texto se ve varias veces.
\end{definicion}

\begin{proposicion}\label{prop:factible}
Sea $E_{\max}>0$. Existe una mezcla con $e_i\le E_{\max}$ para todo $i$ si y solo si $D\le E_{\max}\sum_{i}U_i$.
\end{proposicion}

\begin{proof}
Si $e_i\le E_{\max}$ para todo $i$, entonces $w_iD\le E_{\max}U_i$; sumando sobre $i$ y usando $\sum_iw_i=1$, $D\le E_{\max}\sum_iU_i$. Recíprocamente, si se cumple la desigualdad, los pesos proporcionales $w_i=U_i/\sum_jU_j$ dan $e_i=D/\sum_jU_j\le E_{\max}$ para todo $i$.
\end{proof}

\begin{ejemplo}[La mezcla de LLaMA]\label{ej:llama-mezcla}
La tabla~1 de~\cite{llama} da, para el entrenamiento de $1{,}4$ billones de tokens, el peso y las épocas de cada fuente. Despejando $U_i=w_iD/e_i$ se obtienen los tokens únicos implícitos:
\begin{center}\small
\begin{tabular}{@{}lrrrr@{}}
\toprule
Fuente & $w_i$ & $e_i$ & $w_iD$ ($10^{9}$) & $U_i=w_iD/e_i$ ($10^{9}$) \\
\midrule
CommonCrawl & $0{,}670$ & $1{,}10$ & $938$ & $853$ \\
C4 & $0{,}150$ & $1{,}06$ & $210$ & $198$ \\
GitHub & $0{,}045$ & $0{,}64$ & $63$ & $98$ \\
Wikipedia & $0{,}045$ & $2{,}45$ & $63$ & $26$ \\
Libros & $0{,}045$ & $2{,}23$ & $63$ & $28$ \\
ArXiv & $0{,}025$ & $1{,}06$ & $35$ & $33$ \\
StackExchange & $0{,}020$ & $1{,}03$ & $28$ & $27$ \\
\bottomrule
\end{tabular}
\end{center}
Wikipedia y los libros, fuentes pequeñas y valiosas, se recorren más de dos veces; parte de GitHub no llega a verse. El \cref{tema:5} cuantifica cuánto vale un token repetido: hasta unas cuatro épocas, casi lo mismo que uno nuevo~\cite{muennighoff}.
\end{ejemplo}

\section{La escala de los datos}

La \cref{tab:tokens-param} recoge la relación entre tokens de entrenamiento y parámetros de varios modelos con informe público. El cambio entre 2020 y 2023 se explica en los temas~\ref{tema:4} y~\ref{tema:5}.

\begin{table}[h]
\centering\small
\caption{Parámetros, tokens de entrenamiento y tokens por parámetro.}\label{tab:tokens-param}
\begin{tabular}{@{}lrrr@{}}
\toprule
Modelo (referencia) & $N$ & $D$ & $D/N$ \\
\midrule
GPT-3~\cite{brown} & $1{,}75\cdot 10^{11}$ & $3\cdot 10^{11}$ & $1{,}7$ \\
Gopher~\cite{rae} & $2{,}8\cdot 10^{11}$ & $3\cdot 10^{11}$ & $1{,}1$ \\
Chinchilla~\cite{hoffmann} & $7\cdot 10^{10}$ & $1{,}4\cdot 10^{12}$ & $20$ \\
LLaMA 7B~\cite{llama} & $6{,}7\cdot 10^{9}$ & $10^{12}$ & $149$ \\
Llama~2 7B~\cite{llama2} & $6{,}74\cdot 10^{9}$ & $2\cdot 10^{12}$ & $297$ \\
Llama~3 405B~\cite{llama3} & $4{,}05\cdot 10^{11}$ & $1{,}56\cdot 10^{13}$ & $38{,}5$ \\
\bottomrule
\end{tabular}
\end{table}

\begin{lecturas}
\item Stanford CS336~\cite{cs336a4}, práctica~4 «Data», §2.6 «Quality Rules», §2.7 «Quality Classifier» y §4 «Filter data for language modeling».
\item Stanford CS324~\cite{cs324}, clase «Data».
\item Wenzek \emph{et al.}~\cite{wenzek}, CCNet: filtro por perplejidad con modelos de $n$-gramas.
\item Penedo \emph{et al.}~\cite{refinedweb}, §3, y~\cite{fineweb}, §3 y §4: procesos completos sobre Common Crawl.
\item Touvron \emph{et al.}~\cite{llama}, §2.1 «Pre-training Data», tabla~1.
\end{lecturas}

\section*{Ejercicios del tema 1}
\addcontentsline{toc}{section}{Ejercicios del tema 1}

\begin{ejercicio}{1}{eje:1-sesgo}
Un equipo entrena un clasificador de calidad con páginas citadas desde la Wikipedia en inglés como ejemplos positivos y lo aplica a un rastreo multilingüe. Indique dos tipos de texto que el filtro tenderá a descartar injustificadamente y una consecuencia para el modelo resultante. ¿Qué medida reduciría el problema?
\end{ejercicio}

\begin{ejercicio}{1}{eje:1-gopher}
Aplique los filtros de la \cref{def:gopher} a los cuatro documentos siguientes y decida cuáles pasan.
\begin{center}\small
\begin{tabular}{@{}lrrrrr@{}}
\toprule
Doc. & Palabras & Caracteres en palabras & Líneas & Líneas con «...» & Palabras con letra \\
\midrule
A & $40$ & $200$ & $4$ & $0$ & $38$ \\
B & $1200$ & $5640$ & $60$ & $24$ & $1170$ \\
C & $800$ & $4160$ & $50$ & $6$ & $720$ \\
D & $3000$ & $13\,500$ & $100$ & $5$ & $2250$ \\
\bottomrule
\end{tabular}
\end{center}
\fuente{adaptado de Stanford CS336~\cite{cs336a4}, práctica~4, §2.6, problema \texttt{gopher\_quality\_filters} (se sustituye la implementación por el cálculo).}
\end{ejercicio}

\begin{ejercicio}{2}{eje:1-mezcla}
Se entrena con $D=10^{12}$ tokens una mezcla de cuatro fuentes: web ($8\cdot 10^{11}$ tokens únicos, peso $0{,}60$), código ($2\cdot 10^{11}$, $0{,}25$), libros ($5\cdot 10^{10}$, $0{,}10$) y ciencia ($10^{10}$, $0{,}05$). Calcule las épocas de cada fuente. ¿Cuál preocupa, si se acepta un máximo de unas $4$ épocas? Proponga un ajuste de pesos.
\end{ejercicio}

\begin{ejercicio}{2}{eje:1-ppl}
Un modelo de unigramas de referencia asigna $q(\text{de})=0{,}4$, $q(\text{la})=0{,}3$, $q(\text{casa})=0{,}2$ y $q(\text{oferta})=0{,}1$. Calcule la perplejidad de los documentos A~=~«la casa de la», B~=~«oferta oferta oferta casa» y C~=~«de de de de». ¿Cuáles conserva un filtro con umbral $\PPL\le 4$? Compruebe la \cref{prop:ppl-kl} para el documento~A.
\end{ejercicio}

\begin{ejercicio}{2}{eje:1-fermi}
Un proceso de filtrado procesa $2500$ archivos WET de Common Crawl en $1{,}5$ horas en una máquina. Supóngase que un volcado completo tiene $100\,000$ archivos WET. (a) ¿Cuánto tarda en filtrarse un volcado en esa máquina? (b) ¿Cuántas máquinas hacen falta para filtrarlo en $2$ horas? (c) ¿Cuántas horas-máquina requieren los $96$ volcados de FineWeb?
\fuente{adaptado de Stanford CS336~\cite{cs336a4}, práctica~4, §4, problema \texttt{filter\_data}, apartado~(b) (con tiempos supuestos).}
\end{ejercicio}

\begin{ejercicio}{3}{eje:1-minmax}
Con la notación de la \cref{def:mezcla}, demuestre que, para $D$ fijo, la mezcla que minimiza el máximo de las épocas, $\max_ie_i$, es la proporcional, $w_i=U_i/\sum_jU_j$, y que es la única. Calcule ese mínimo para la mezcla de LLaMA del \cref{ej:llama-mezcla} y compárelo con las épocas de Wikipedia elegidas por sus autores.
\end{ejercicio}
